% ============================================================================
%  makroi.tex — notation macros. One definition per concept so the whole book
%  stays consistent with the professor's notation (Gx - h <= 0, theta(mu), v).
% ============================================================================

\providecommand{\coloneqq}{\mathrel{\vcentcolon=}}

\newcommand{\R}{\mathbb{R}}                      % real numbers
\newcommand{\Rn}[1]{\mathbb{R}^{#1}}
\newcommand{\Sym}[1]{\mathbb{S}^{#1}}            % symmetric matrices
\newcommand{\trans}{^{\mathsf{T}}}               % transpose
\newcommand{\grad}{\nabla}
\newcommand{\Lag}{\mathcal{L}}                   % Lagrangian
\newcommand{\dual}{\theta}                       % dual function theta(mu)
\newcommand{\feas}{\mathcal{F}}                  % feasible set
\newcommand{\popt}{p^{\star}}                    % optimal primal value
\newcommand{\dopt}{d^{\star}}                    % optimal dual value
\newcommand{\xopt}{x^{\star}}
\newcommand{\muopt}{\mu^{\star}}
\newcommand{\psd}{\succeq}                       % positive semidefinite
\newcommand{\pd}{\succ}                          % positive definite
\newcommand{\nsd}{\preceq}
\newcommand{\nd}{\prec}
\newcommand{\norm}[1]{\left\lVert #1 \right\rVert}
\newcommand{\abs}[1]{\left\lvert #1 \right\rvert}
\newcommand{\inner}[2]{\left\langle #1,\, #2 \right\rangle}
\newcommand{\diag}{\operatorname{diag}}
\newcommand{\eig}{\operatorname{eig}}
\newcommand{\rank}{\operatorname{rank}}
\newcommand{\tr}{\operatorname{tr}}
\newcommand{\Real}{\operatorname{Re}}
\newcommand{\dom}{\operatorname{dom}}
\newcommand{\dd}{\,\mathrm{d}}               % differential in integrals
\newcommand{\subjto}{\text{uz ograničenja}}

% Reference to a page in Part 0 where a symbol is introduced.
\newcommand{\vidi}[1]{\textcolor{tihasiva}{\footnotesize(vidi str.~\pageref{#1})}}
